\section{Related work}
\label{sec:related}

The problem sits where three bodies of work meet: routing and cascading
across models, control of reasoning effort at inference time, and
measurement of the energy inference consumes. Table~\ref{tab:prior-work}
places the closest systems side by side and
Section~\ref{sec:related-gap} reads it.

\subsection{Routing and cascading across models}
\label{sec:related-routing}

This paper is concerned with routing that sends whole queries to separate
models, each of which could answer alone, rather than with routing tokens
between expert sub-networks inside one model. Moslem and Kelleher survey that
literature and divide it into routing, which commits before generation, and
cascading, which tries models in sequence \cite{moslem2026routingsurvey}.

\emph{Cascades.} FrugalGPT sends a query through models of increasing price
and stops when a learned scorer accepts the answer, reporting up to 98\%
lower cost than the best individual model \cite{chen2023frugalgpt}. AutoMix
replaces the external scorer with few-shot self-verification by the small
model and decides whether to escalate with a partially observable Markov
decision process, reporting cost reductions above 50\%
\cite{aggarwal2024automix}. AutoMix is the closest published idea to the
question asked here, because it treats escalation as a sequential decision;
its decision state, however, is the model's confidence in its own answer
rather than the cost already spent, and it is an automated policy rather than
a model of what a user does. Both cascades share a structural cost: the cheap
model must generate before the system can decide, so every escalated query
pays for at least one discarded answer.

\emph{Predictive routers.} Hybrid LLM estimates from the query alone whether
a small model's answer will match a large model's, and exposes a threshold
that trades quality for cost, reporting 40\% fewer large-model calls at no
loss in quality \cite{ding2024hybridllm}. RouteLLM learns a strong-versus-weak
decision from human preference data and reports more than a twofold cost
reduction \cite{ong2024routellm}. Later work trains a query encoder
with dual contrastive losses \cite{chen2024routerdc}, routes to models unseen
during training with an excess-risk bound \cite{jitkrittum2025universal},
aligns a router with user-defined preferences so models can be added without
retraining \cite{tran2025archrouter}, and interleaves a router's own
deliberation with calls to other models over several rounds
\cite{zhang2025routerr1}. Published benchmarks make routers comparable, among
them LLMRouterBench, whose released outcomes and prices this study reuses
\cite{li2026llmrouterbench}. Across this work, cost is denominated in API
price, model calls or tokens, and the comparison points are static, random,
oracle or automated policies.

\subsection{Reasoning effort and test-time compute}
\label{sec:related-reasoning}

Chain-of-thought prompting improves multi-step reasoning by spending
generated tokens on intermediate steps \cite{wei2022cot}, reinforcement
learning now trains models that produce long traces unprompted
\cite{deepseek2025r1}, and Qwen3 builds the control into the model, serving
thinking and non-thinking behaviour from one set of weights \cite{qwen3}.
Long traces are not always useful: Sui
et al.\ document reasoning models spending many tokens on easy questions
\cite{sui2025stopoverthinking}, and Alomrani et al.\ organise adaptive
test-time compute into controllability under a fixed budget and adaptiveness,
where the budget depends on the input \cite{alomrani2025reasoningbudget}.
Selecting a reasoning mode per prompt is an instance of the second tier.

A recent group of papers routes over whether and how much to think, usually
within one model. SynapseRoute sends medical questions to the thinking or
non-thinking mode of a dual-state model, reporting 36.8\% lower inference
time and that about 58\% of its questions need no reasoning
\cite{zhang2025synapseroute}. A lightweight router head can make the same
choice for ranking tasks \cite{guo2026thinkwhenneeded}, while DART is
training-free, reading difficulty from the agreement of two cheap drafts and
cutting thinking tokens by 32--73\% \cite{lee2026dart}. Learning
When to Think moves the decision inside the model, which emits NoThink, Short
or Long as its first token \cite{kassenaar2026learningwhentothink}, and
HRBench finds that the preferred switching strategy depends on model scale
and task domain \cite{ning2026hrbench}. Route to Reason is the closest prior
work to the grid studied here, selecting jointly among seven models and four
reasoning strategies under a budget constraint \cite{pan2026routetoreason};
joint selection of model and reasoning effort is therefore not, by itself, a
new idea, and we do not claim it as one.

All of this work depends on an assumption that is rarely tested: that a
model's thinking and non-thinking modes are genuinely different operating
points. Wang et al.\ found only partial separation in hybrid-thinking models,
with reasoning behaviour leaking into non-thinking outputs
\cite{wang2025demystifying}, which is why we test separation on our own
checkpoints before labelling (Section~\ref{sec:setup-gates}).

\subsection{Energy, carbon and energy-aware routing}
\label{sec:related-energy}

Strubell et al.\ brought the cost of training to the field's attention
\cite{strubell2019energy} and Schwartz et al.\ argued for efficiency as a
first-class result \cite{schwartz2020greenai}. Lacoste et al.\ convert
compute into emissions with regional grid intensity
\cite{lacoste2019quantifying}, and Patterson et al.\ show how strongly
training emissions depend on hardware and location \cite{patterson2021carbon}.
For inference, Luccioni et al.\ measured energy and carbon per 1{,}000
queries across task types and found wide variation
\cite{luccioni2024powerhungry}. Tuning GPU frequency and batch size
\cite{xu2025camel} or managing key--value cache memory \cite{kwon2023vllm}
makes each configuration cheaper to run without deciding which to run, and is
therefore orthogonal to routing.

Three routers treat energy or latency as part of the objective. GreenServ
routes among 16 open models, measures GPU energy on an A100 and reports a
31\% energy reduction with 22\% higher accuracy than random routing, but does
not subtract the router's own energy \cite{ziller2026greenserv}. EcoThink
sends each query to a retrieval or a reasoning path with a distilled
classifier and reports a 40.4\% energy reduction and 1.32\,gCO$_2$e per query
against 4.52 for GPT-4o, with energy partly estimated from thermal design
power and the router uncharged \cite{li2026ecothink}. Yang et al.\ route
among edge experts on RTX~4090 GPUs, optimising quality of service and
latency without measuring energy \cite{yang2025qosrouting}. Energy-aware
routing and carbon reporting are therefore established, and neither is a
contribution of this paper; what is missing is a router over reasoning mode
whose own measured energy is netted out.

Open-ended prompts cannot be graded by string match, so studies that include
them rely on model judges. Zheng et al.\ found a strong judge agreeing with
human preferences about as often as humans agree with each other, while
documenting position, verbosity and self-enhancement biases
\cite{zheng2023judging}. Those biases motivate the out-of-family judge and
two-sided validation of Section~\ref{sec:setup-judge}.

\subsection{Comparison with related work and positioning}
\label{sec:related-gap}

\begin{table*}[ht]
\centering
\footnotesize
\caption{The closest prior systems compared with this study. \emph{Routes
over} is what the system selects, \emph{decides} is when it commits, and the
last four columns record what each work reports. The saving is the headline
figure its authors quote against their own baselines, which differ from
system to system and from ours. A dash means the quantity does not apply or
is not part of the work's reported results.}
\label{tab:prior-work}
\renewcommand{\arraystretch}{1.15}
\setlength{\tabcolsep}{4pt}
\begin{tabularx}{\textwidth}{@{}l >{\raggedright\arraybackslash}p{2.35cm} >{\raggedright\arraybackslash}p{2.35cm} >{\raggedright\arraybackslash}X >{\raggedright\arraybackslash}p{1.85cm} >{\raggedright\arraybackslash}p{2.75cm} >{\raggedright\arraybackslash}p{1.6cm}@{}}
\toprule
System & Routes over & Decides & Efficiency measure and reported saving & Energy & Router's own cost & Manual escalation \\
\midrule
FrugalGPT \cite{chen2023frugalgpt} & Model & After each answer (cascade) & Up to 98\% lower cost & --- & --- & --- \\
AutoMix \cite{aggarwal2024automix} & Model & After self-verification (cascade) & Over 50\% lower cost & --- & --- & --- \\
Hybrid LLM \cite{ding2024hybridllm} & Model, small or large & Before generation & 40\% fewer large-model calls & --- & --- & --- \\
RouteLLM \cite{ong2024routellm} & Model, strong or weak & Before generation & Over twofold lower cost & --- & --- & --- \\
Router-R1 \cite{zhang2025routerr1} & Model & Over several rounds & Parameter-weighted output tokens & --- & --- & --- \\
SynapseRoute \cite{zhang2025synapseroute} & Reasoning mode & Before generation & 36.8\% less inference time & --- & --- & --- \\
DART \cite{lee2026dart} & Thinking budget & After two short drafts & 32--73\% fewer thinking tokens & Not measured & --- & --- \\
Learning When to Think \cite{kassenaar2026learningwhentothink} & Thinking level & In the model, at the first token & Response length & --- & --- & --- \\
Route to Reason \cite{pan2026routetoreason} & Model and reasoning strategy & Before generation & Over 60\% fewer tokens & --- & --- & --- \\
GreenServ \cite{ziller2026greenserv} & Model & Before generation & 31\% less energy & Measured, A100 & Latency reported, energy not netted & --- \\
EcoThink \cite{li2026ecothink} & Retrieval or reasoning path & Before generation & 40.4\% less energy & Partly estimated & Not netted & --- \\
QoS edge routing \cite{yang2025qosrouting} & Edge expert & Before generation & QoS score under a latency bound & Not measured & Latency reported & --- \\
\midrule
This study & Model size and reasoning mode & Before generation & Energy per solved prompt; 61\% less than the sticky default, router included & Measured, board power & Netted: tokens, latency, energy & Yes, four policies \\
\bottomrule
\end{tabularx}
\end{table*}

Read along its rows, Table~\ref{tab:prior-work} shows three differences. The
first is the axis. Most of these systems route over models alone or over
reasoning mode alone; Route to Reason is the closest to the grid studied here,
selecting jointly from seven models and four reasoning strategies
\cite{pan2026routetoreason}. We therefore do not claim joint selection as
new. The second is the measure of efficiency. Savings
are quoted in price, model calls, tokens or time; only GreenServ meters energy
on hardware and EcoThink partly estimates it, and both report savings gross of
the router, as the table's last-but-one column records. The third is the
comparator. Every baseline in these studies is a static policy, random
assignment, an oracle or an automated cascade, and none is the manual
escalation, with the prompt resent at each step, that this study prices. The
quoted savings are not comparable with ours: they are measured against
different baselines, on different pools and hardware, and in different units,
which is why we place them in a table and do not set them against our own
figures.

Four questions remain open in the literature we located. First, the
escalation tax is unmeasured: no study formalises manual escalation, in which
each failed attempt is billed in full and the prompt is resent. Second,
router energy is not netted out; we found no study that subtracts the
selector's own measured energy, or asks when a language-model router stops
paying for itself against a linear classifier. Third, the reasoning-mode axis
has no measured energy cost: routers over reasoning mode report tokens and
latency, and routers that measure energy do not route over mode. Fourth,
energy is not normalised by success; energy per \emph{solved} query, which
charges a cheap failure for the retry it causes, was not found.
